\chapter{Describing Wisconsin School Board Elections}
\label{chap:wisco}
\chapterstyle{deposit}




\section{The Wisconsin Case}

As noted in Chapter \ref{chap:intro}, this dissertation focuses on the case of a 
single state -- Wisconsin. Wisconsin is located in the Midwest and ranks 
23rd in terms of total area, 20th in terms of population size with 5.7 million residents, 
and has over 420 local school districts. Recently, Wisconsin has become a crucial swing state 
exhibiting a close split between Republican and Democratic voters in presidential 
politics. I have selected Wisconsin for a number of reasons. Wisconsin, like many 
other states in the Midwest is a state that grants much of the power and authority 
in school district decision making to local school boards -- a so-called ``local 
control'' state. This is also why Wisconsin, in contrast to Southern states, 
has such a large number of school districts with relatively small sizes. 
Prior studies of school boards like those of \citet{RefWorks:311,RefWorks:214} have 
taken place in states with larger school districts and more centralized state 
control. Thus, Wisconsin serves as a good exemplar for other ``local control'' 
jurisdictions where school boards have more authority relative to the state 
government. Another advantage of selecting Wisconsin is the comparative advantage
I have had in gathering the necessary data to study a substantial proportion of the 
school board races within the state.\footnote{For details on this, see the 
Appendix \ref{app:data}.} A final advantage is the ability 
to test theories of the linkage between school board and state level politics thanks 
to the dramatic shock to education policy that accompanied the election of Wisconsin 
Governor Scott Walker in 2010. Taken together, these advantages mean that the in-depth 
study of school board elections across the entire K-12 educational system in Wisconsin 
can provide a new perspective on theories of candidate emergence in local elections, 
voter turnout in non-partisan off-cycle local races, and federalist linkages between 
state and local governments. 

In this chapter I address the unique case of Wisconsin by providing a descriptive 
overview of the legal and the electoral landscape for school board races within 
the state. I start by describing the ways in which the laws governing 
school board races vary from the more familiar legislative races. Next, I describe the data 
I have collected on school board election results in detail and provide descriptive 
statistics of the frequency of incumbency challenges and open races, voter turnout, 
and the socio-demographic diversity of the constituencies served by Wisconsin school 
boards. This overview serves as the groundwork for the three subsequent chapters 
which examine how well theories of political engagement explain candidate emergence, 
voter turnout, and policy shifts in school board elections. 

\section{Background and Legal Basis of Wisconsin School District Elections}
\label{sec:wiscolaw}
This project represents the first attempt to comprehensively study the patterns 
in school board elections across an entire state over multiple years. Before 
examining the patterns that emerge both within and between school districts, 
it is important to first explore the data available on school board elections. 
Wisconsin, like many other Midwest states and unlike states in the South, is a 
local control state. This means that Wisconsin school districts have revenue 
authority independent of municipalities and counties, and are governed by 
a non-partisan elected board independent of counties and other local government 
bodies. The independence of Wisconsin school districts is established in the 
state constitution and through state law (WI Const. art. X, \S 3 and Wis. Stat 120.10). 

This independence gives Wisconsin school districts considerable power and authority. 
Importantly, school boards administer their own elections and retain their own 
election records. This represents a serious barrier to studies of school board 
elections due to the burden associated with collecting and organizing school board 
elections from over 420 administrative entities. This is why previous studies of 
school boards have tended to focus on candidates, random samples of school boards, 
or on states with county-wide school boards such as North Carolina or Georgia. 

Wisconsin school board elections are held consistently on the same dates across 
the state, during the spring non-partisan election cycle. The spring primary election is the 
third Tuesday in February, and the spring general election is the first Tuesday 
in April (Wis. Stat \S 120.06). Wisconsin school boards must have between 3 and 9 
members with some exceptions made for boards of 11 members (Wis. Stat \S 120.01). 
Board member terms are not consistent across school boards, though 3 year terms are 
the most common with limited exceptions for terms of 1 or 2 years (Wis. Stat \S 120.06 (3)). 
Board members are elected by a plurality and must be residents of the district to 
be eligible for office (Wis. Stat \S 120.06 (2) (a)). All candidate verification, 
ballot preparation, and election announcement is the responsibility of the school district 
clerk (Wis. Stat \S 120.06 (8)). Primary elections occur if there are more than 
twice as many candidates as there are members to be elected (Wis. Stat \S 120.06 (7) (b)). 
School boards may choose to apportion their seats in dramatically different ways 
ranging from a board of 3-9 members each representing the entire district, to a board 
of 3-11 members each representing either a specific region of the school district or 
the entire district. The only restriction on the way seats are allocated is that the 
plan is proposed with a petition by a threshold of voters and passed by the board 
at an annual meeting or through an election (Wis. Stat \S 120.02 (2)). 


\begin{table}[h!]
\begin{center}
\caption{Wisconsin Election Dates}
\label{tab:elecdates}
\begin{tabular}{| p{.3\textwidth} | p{.6\textwidth} |}
\hline
\textbf{Election}  & \textbf{Date}  \\ \hline
Spring Primary  & Third Tuesday in February  \\ \hline
Spring General &  First Tuesday in April  \\ \hline
Fall Primary & Second Tuesday in August \\ \hline
Fall General & First Tuesday in November\\ \hline
\end{tabular}
\end{center}
\end{table}

My study distinguishes itself from many other studies of school board elections through 
its use of official election results provided directly by school boards to determine 
the results of the elections. These election records came in many forms. The official and preferred 
records included the Statement of the Board of Canvass and the Certificate of the 
Board of Canvass. These two documents are legally mandated for each school board 
election and certify the winner for each school board seat, as well as the votes 
cast for each candidate. In most cases the records were provided electronically, 
though some were also received as hard copy in the mail. Full details on the data 
collection process can be found in Appendix \ref{app:data}. 
These records were then transcribed from the various records provided by school 
districts into spreadsheets that combined together form the key database for this 
analysis. 

Using the names, vote totals, and winners from each school board election allow 
for the construction of measures of the number of candidates in each election 
cycle, the number of contested seats, the level of voter participation, and other 
metrics of electoral activity at the school board level. Importantly, because the 
data was collected for several years in each district, time-variant trends can be 
explored. 

In Wisconsin, school boards are the official holder of all school district elections. 
This means that official school board election records are maintained solely by 
school districts themselves and gathering school board election results in Wisconsin 
requires requesting records from individual school districts.\footnote{See Appendix 
\ref{app:data} and the Appendix \ref{app:log} for details on the request process.} 
An additional challenge is that though school boards are required by law to retain 
these records, the record keeping is often incomplete, unclear, or confusing (Wis. 
Stat \S 120.06 (8) (g)).\footnote{While a few enterprising county clerks have 
collected school board election results and made them available, these can be incomplete 
by only including the precincts that are within the county. They also are not official. 
\citep[personal communication with][]{clerk1}.}In order to conduct this study then, a records 
request was made to all 424 public school districts in the state of Wisconsin. 
This request was e-mailed to the school district administrator in each of Wisconsin's 
424 school districts in January of 2013. Three follow up waves were conducted throughout 
2013, and at the end of the collection period, 312 
districts responded with records. The median district provided 
11 
years of records.\footnote{While each district was requested to provide the maximum amount of records legally 
required to be retained, ten years, many districts were unable to locate complete 
records going this far back.} Overall, records were provided on 4,181 school 
board races involving 5,863 candidates in the state. For each election 
the votes received by each candidate was recorded, as well as whether the candidate 
won, was an incumbent, was a repeat candidate in the school district, or was a 
minor candidate. For each race, whether the candidate was running for district 
wide or regional seat was also recorded. 

\begin{knitrout}
\definecolor{shadecolor}{rgb}{0.969, 0.969, 0.969}\color{fgcolor}\begin{figure}[]

\includegraphics[width=\maxwidth]{../../phd_figs/chap3-yeargraph} \caption[Number of Districts Providing Records by Year]{Number of Districts Providing Records by Year\label{fig:yeargraph}}
\end{figure}


\end{knitrout}



\begin{knitrout}
\definecolor{shadecolor}{rgb}{0.969, 0.969, 0.969}\color{fgcolor}\begin{figure}[]

\includegraphics[width=\maxwidth]{../../phd_figs/chap3-racesgraph} \caption[Number of Races Provided by Districts]{Number of Races Provided by Districts. Red line indicates the median.\label{fig:racesgraph}}
\end{figure}


\end{knitrout}





The good news is that the most crucial period of the sample, from 2007 to 2012, is 
the period with the most complete set of records. The mean district provided records on 
13.4 school 
board races. 283 districts had 
a race in each year between 2007 and 2012, while 29 
did not. I distinguish these cases from others by verifying the absence of a ballot 
for school board, as opposed to the absence of records. 

\section{The Diversity of Wisconsin School districts}

The authority given to school districts in Wisconsin to determine their own electoral 
rules and structure make Wisconsin a useful case in exploring the variation in 
school board electoral participation across contexts. In addition to the legal 
variation in electoral rules across school districts, Wisconsin school districts 
also vary widely in size, wealth, performance, and student and household demographics. In order to 
explore this diversity and to demonstrate the representativeness of the school 
districts that responded to the request for records, I now explore several aspects 
of Wisconsin school districts. 

\subsection{Comparing the Sample to the Population}











































